\documentclass{discussion}

\usepackage{amsmath, amsthm, amssymb}
% \usepackage{xcolor, hyperref, cleveref}
\usepackage{mathtools}
\usepackage{derivative}
\usepackage{tabularx}
\usepackage{graphicx}
\usepackage{tasks}
\usepackage{mdframed}
\usepackage{interval}
\usepackage{commands}

\settasks{
    label-width = 13.99168pt,
    item-indent = 17.6413pt,
    after-item-skip = 80pt,
}

\discussion{5}
\author{Sections B04 and B06}
\date{Fall 2024}

\begin{document}

\begin{question}
    Estimate the error if \(\displaystyle P_3(x) = 1 + x + \frac{x^2}{2} + \frac{x^3}{6}\) is used to estimate the value of \(f(x) = \eu^{x}\) by filling in the blanks below.
    \begin{enumerate}
        \item By Taylor's theorem, for each \(x\), the error is
        \begin{equation*}
            R_3(x) = f(x) - P_3(x) = \frac{f^{(4)}(c)}{4!} x^{4} = \rule{2cm}{0.15mm} \; x^{4}
        \end{equation*}
        for some \(c\) in between \(0\) and \(x\).

        \item Noting that \(\abs{c} \leq \abs{x}\), the error is bounded by (fill in a number)
        \begin{equation*}
            \abs{R_3(x)} \leq \rule{1cm}{0.15mm} \; \eu^{\abs{x}} \abs{x}^4.
        \end{equation*}

        \item If we were to consider all \(x\) in \(\ointerval{-1}{1}\), an upper bound for the error is (fill in a number)
        \begin{equation*}
            \abs{R_3(x)} \leq \rule{1cm}{0.15mm}
        \end{equation*}

        \item In general, if we were to consider all \(x\) in \(\ointerval{-a}{a}\), an upper bound for the error is
        \begin{equation*}
            \abs{R_3(x)} \leq \rule{3cm}{0.15mm}
        \end{equation*}
        depending on the value of \(a\).

        If we want to find an interval \(\ointerval[scaled]{-\frac{1}{N}}{\frac{1}{N}}\) such that the error is less than \(0.001\), we need to solve the inequality
        \begin{equation*}
            \abs{R_3(x)} \leq \rule{3cm}{0.15mm} < 0.001,
        \end{equation*}
        so one possible value of \(N\) is \(\rule{1cm}{0.15mm}\).
    \end{enumerate}
\end{question}

\begin{question}
    Calculate an upper bound for the error if the following \(P_n(x)\) is used to estimate \(f(x)\) on \(\ointerval{-0.1}{0.1}\).
    Moreover, find an interval \(\ointerval{-a}{a}\) such that the error is less than \(0.001\).
    \begin{tasks}(2)
        \task \(\displaystyle P_1(x) = 1 + \frac{x}{2}\) for \(f(x) = \sqrt{1 + x}\)
        \task \(\displaystyle P_4(x) = x - \frac{x^3}{6}\) for \(f(x) = \sin(x)\) \\
        \tiny (this \(P_4\) is indeed a fourth-degree approximation of \(\sin(x)\))
        % \newline (note that \(P_4\) is indeed a fourth-degree approximation of \(\sin(x)\))
        % \task \(\displaystyle P(x) = 1 + x + \frac{x^2}{2} + \frac{x^3}{6} + \frac{x^4}{24}\) for \(f(x) = \eu^{x}\)
        \task \(\displaystyle P_2(x) = x + x^{2}\) for \(\displaystyle f(x) = \sin(x) \eu^{x}\)
        \task \(\displaystyle P_{3}(x) = \frac{x^3}{3}\) for \(\displaystyle f(x) = \int_{0}^{x} \sin(t^2) \odif{t}\)
    \end{tasks}
\end{question}

% \section{Determine interval with given error bound}

\end{document}